\section{Modelling Decisions}
\label{sec:decisions}

Seven decisions shape both metamodels. Each is stated with the observation that
forced it, because in every case the alternative had already been tried in the
estate's hand-maintained configuration.

\subsection{Three Model Stages, Two Metamodels}

The pipeline has three stages: an authored model, a resolved model, and the
generated text. Only the first two are given metamodels. The generated files are
text in a target syntax that already has a schema of its own, so giving them a
third metamodel would restate that schema without adding a constraint the
project could check.

The three stages are deliberately \emph{not} described as metalevels. No stage
is a type model of the stage below it, so calling the arrangement a metamodel
hierarchy would claim an \texttt{instanceOf} relation that does not hold. They
are stages of one transformation pipeline, and the word metamodel is reserved
for a language definition.

\subsection{The Layer Rule}

The rule that makes the middle stage worth having: \textbf{the authored model
contains no mechanisms, and the generated text contains no decisions.} Every
value is assignable to exactly one stage by two questions. Must a human author
it? And does it record a decision, or merely serialise one?

The counter-experiment is on record. An earlier two-stage arrangement, with
resolution private to the generator, produced three mutually incompatible
documents all claiming the same version identifier, because a two-stage
vocabulary could not say which of the three was the authoring shape. The middle
model exists so that a decision has one place to live and one name.

\subsection{Contention Decides Authority}

A value is decided by the platform if and only if it must be unique across the
estate or draws on a shared finite resource; every other value is declared by
the application and carried through untouched. One question replaces a
per-field negotiation.

The rule required one refinement that is worth reporting, because the naive
reading is wrong. Contention decides who \emph{arbitrates}, not who
\emph{authors}. Memory and CPU are draws on a finite pool, so a
who-authors reading would forbid an application from stating them, which leaves
every process with no reservation at all. The application therefore states its
requirement and the platform arbitrates it against what the machines publish.
The accepted cost is named rather than hidden: nothing today stops an author
overstating the requirement, because the arbitration is an eligibility test
against one machine rather than a sum across the estate.

\subsection{Derivation Is Total, and There Is No Override}

Every value the platform decides is a function of the authored model and the
platform's own facts. A derived value has exactly one declaring site, so the
metamodel carries no override mechanism and no exception vocabulary.

This decision was taken against evidence that initially argued for a hatch: one
process carried a hand-set deadline three times smaller than its neighbours,
justified by the observation that it starts in seconds rather than minutes. That
is not a value only its owner could know; it is evidence that one rule over the
startup budget was wrong for a whole class of process. The correct response is a
rule that reads an input the process already declares. Keeping the hatch would
have hidden a wrong derivation behind a per-process reason, where removing it
makes the same error visible as a wrong result for the entire class.

\subsection{Values Name Model Concepts, Never Target Fields}

No attribute of the source metamodel, and no attribute of the target metamodel,
is named after a field of the generated syntax. A process declares whether it
may be taken down during a cutover, not a rollout strategy; it declares the
paths it must write, not a security context. The target metamodel records the
platform's decisions in the same vocabulary, and the generators are the only
place a target field name is spelled.

This is what keeps the target metamodel a model rather than an object tree with
extra steps, and it is checked directly: no key of either metamodel may be a
field name of the target syntax.

\subsection{Closed Vocabularies}

Wherever an attribute admits a fixed set of values, the set is closed in the
metamodel rather than typed as a string. The estate supplies the argument: one
existing field was typed as a free-form string, every application wrote the same
value into it, and not one of those values was correct. A field that is wrong in
every observed case at no cost is a comment, not a declaration, and closing the
vocabulary is what turns it back into one.

\subsection{Structure Carries Atomicity; the Graph Carries Order}

Two groupings are needed, and deriving one from the other was tried and
rejected. What must be released together is \emph{structural}: the processes of
one application switch together, or none of them switches. What must be applied
before what is \emph{derived} from the dependency graph. Deriving atomicity from
the graph instead takes the transitive closure of every edge and turns the whole
estate into one unit, which would make every change estate-wide.
